Peripersonal space (PPS) is often inferred from tactile reaction time during an approaching sound. We tested whether instructed respiratory phase changed this SOA-dependent response differently in breath-coupled virtual reality (Viscereality) and Control. The released cohort comprised 42 participants and 8{,}916 retained audio--tactile hit trials. The planned Condition $\times$ Phase $\times$ SOA mixed-model interaction was not detected, $\chi^{2}(1)=0.024$, $p=.878$; categorical-SOA and partial-session-exclusion sensitivities agreed. Audio--tactile responses were 18.8~ms faster than tactile-only Baseline responses on average, but equal nominal delay labels do not establish matched cue times. None of 168 participant-by-condition-by-phase trial-level sigmoid fits met the planned $R^{2}\geq .70$ threshold. The maximum saturated five-SOA trial-level $R^{2}$ ceiling was 0.430, precluding this threshold for any curve of SOA alone. Boundary, sharpness, and boundary-related questionnaire inference therefore remained unestimated. The full data-supported pipeline, diagnostics, and downloadable results run in a marimo browser companion. These data do not demonstrate the predicted Viscereality-specific respiratory modulation; confidence intervals and source-timing and respiratory-compliance limits constrain interpretation.
